\section[A fundamental exact sequence]{A fundamental exact sequence. Descent by morphisms with
relatively connected fibers}
\label{IX.6}
% original p. 253

\begin{theorem}
\label{IX.6.1}
Let $S$ be the spectrum of an artinian ring $A$ with residue field
$k$, $\overline{k}$ an algebraic closure of~$k$, $X$ an $S$-prescheme,
$X_0=X\otimes_A k$, $\overline{X}_0=X\otimes_A
\overline{k}$, $\overline{a}$ a geometric point
of~$\overline{X}$, $a$ its image in~$X$, $b$ its image in~$S$. Suppose
that $X_0$ is quasi-compact and \emph{geometrically
connected} over~$k$ (N.B. if $X$ is proper over~$S$, this means that
$\H^0(X_0,\cal{O}_{X_0})$ is a \emph{local} artinian ring with residue
field \emph{radicial} over~$k$). Then the sequence of canonical
homomorphisms
$$
e\mto \pi_1(\overline{X}_0,\overline{a})\to
\pi_1(X,a)\to\pi_1(S,b)\to e
$$
is exact, and one has
$$
\pi_1(S,b)\isomfrom \pi_1(k,\overline{k})=\text{Galois group
of~$\overline{k}$ over~$k$.}
$$
\end{theorem}

Since fundamental groups do not change on dividing out by the
nilpotent elements, one may assume $A=k$, which already makes the last
isomorphism evident. Let $k'$ be the separable closure of~$k$
in~$\overline{k}$, and consider $X'=X\otimes_k k'$, and the image $a'$
of~$\overline{a}$ in~$X'$. One has a sequence of canonical
homomorphisms
$$
e\mto \pi_1(\overline{X}_0,\overline{a})\to
\pi_1(X',a')\to\pi_1(S',b')\to e
$$
(where $S'=\Spec(k')$). Finally, one has a canonical homomorphism from
this sequence into the one relative to~$X/k$, thanks to the~diagram
$$
\xymatrix{ S & X \ar[l] & \overline{X}_0 \ar[l] \\ S' \ar[u] &
X'\ar[u] \ar[l] & \overline{X}_0 \ar[l] \ar@{=}[u]. }
$$
On the other hand, one sees that this homomorphism of sequences of
groups is an isomorphism, as follows from~\Ref{IX.4.11}. One is
therefore reduced to proving that the second sequence is exact, \ie one
may assume that $k$ is \emph{perfect}. Let then $k_i$ be the finite
Galois subextensions of~$k$ in~$\overline{k}$, set $X_i=X\otimes_k
k_i$, and let $a_i$ be the image of~$\overline{a}$ in~$X_i$. We leave
to the reader the task of verifying that the natural homomorphism
$$
\pi_1(\overline{X}_0,a)\isomto \varprojlim\pi_1(X_i,a_i)
$$
% original p. 254
is an isomorphism, which simply means that an étale covering
of~$\overline{X}$ comes from an étale covering of some~$X_i$, and that
the latter is essentially unique, up to passing to some~$X_j$,
$j\geq i$. On the other hand, let $\pi_i$ be the Galois group of~$k_i$
over~$k$, \ie the group opposite to the group of $S$-automorphisms
of~$S_i=\Spec(k_i)$. Since the functor $S'\mto X\times_S S'$ from the
étale coverings of~$S$ to the étale coverings of~$X$ is fully faithful
\eqref{IX.3.4}, it follows that $\pi_i$ is also isomorphic to the group
opposite to the groups of $X$-automorphisms of the connected principal
covering $X_i$ of~$X$. It therefore follows from V~\Ref{V.6.13} that
one has an exact sequence
$$
e\mto \pi_1(X_i,a_i)\to \pi_1(X,a)\to\pi_i\to e.
$$
Passing to the projective limit over~$i$ in these exact sequences, one
finds an exact sequence (since one is in the category of groups of
Galois type), which is none other than the sequence considered
in~\Ref{IX.6.1}. This completes the proof.

The translation of the right exactness in~\Ref{IX.6.1} into geometric
language is the following:

\begin{corollary}
\label{IX.6.2}
With the preceding notation, let $X'$ be an étale covering of~$X$, and
let $\overline{X'}_0$ be the corresponding étale covering
of~$\overline{X}_0$. The following conditions are equivalent:
\begin{enumerate}
\item[(i)] There exists an $S'$ étale over~$S$, and an $X$-isomorphism
$X'\isomto X\times_S S'$, ($S'$ is then determined up to unique
isomorphism by virtue of~\Ref{IX.3.4}).
\item[(ii)] $\overline{X}_0'$ is completely decomposed
over~$\overline{X}_0$.

If $X'$ is connected, these conditions are also equivalent to:
\item[(ii~bis)] $\overline{X}_0'$ has a section over~$\overline{X}_0$.
\end{enumerate}
\end{corollary}

(N.B. This last complement is essential; the equivalence of~(i) and
(ii) means only that $\pi_1(S,b)$ is the quotient group
of~$\pi_1(X,a)$ by the closed invariant subgroup
% original p. 255
generated by the image of~$\pi_1(\overline{X}_0,\overline{a})$, and
not by this image itself). Under the preceding conditions, we shall
say that $X'$ is a \emph{geometrically trivial} covering of~$X$.

\begin{remark}
\label{IX.6.3}
In the statement~\Ref{IX.6.1} one cannot replace $\overline{k}$ by an
arbitrary algebraically closed extension of~$k$, even if $k$ is
already assumed algebraically closed. In other words, it is not true
in general that if $X$ is a connected algebraic scheme over an
algebraically closed field $k$, its fundamental group does not change
on replacing $k$ by an algebraically closed extension; this is already
false for example in characteristic $p>0$ for the affine line over~$k$,
because of the phenomena of ``higher ramification'' at the point at
infinity, implying a ``continuous'' structure for the fundamental
group. We shall see however in the next exposé that such phenomena
cannot occur if $X$ is \emph{proper} over~$k$. We shall also show by
transcendental means that the same holds if $k$ is of characteristic
zero.
\end{remark}

\begin{corollary}
\label{IX.6.4}
Suppose that $a$ is localized at an $x\in X$ which is rational
over~$k$ (or more generally, which has a residue field radicial
over~$k$). Then the exact sequence~\Ref{IX.6.1} is split.
\end{corollary}

One may assume $S=\Spec(k)$. If $x$ is rational over~$k$, it
corresponds to a section $S\to X$ of~$X$ over~$S$, transforming $b$
into~$a$, and defining a homomorphism $\pi_1(S,b)\to \pi_1(X,a)$ which
is the desired splitting. If $\kres(x)$ is radicial over~$k$, one
reduces to the preceding case by making the base extension
$\Spec(\kres(x))\to\Spec(k)$.

\begin{theorem}
\label{IX.6.5}
Let $f\colon X\to S$ be a proper and surjective morphism of finite
presentation, with geometrically connected fibers, $X'$ a prescheme of
finite presentation and proper over~$X$, $s$ a point of~$S$, $F=X_s$
the fiber of~$X$ at~$s$, and $F_1'$ a connected component of the fiber
$F'=X'_s$ of~$X'$ at~$s$. In order that there exist an open
neighborhood $X_1'$ of~$F_1'$ in~$X'$, an étale $S$-scheme $S_1'$ and
an $X$-isomorphism $X_1'\isomto S_1'\times_S X$, it is necessary and
sufficient that $X'$ be étale over~$X$ at the points of~$F_1'$, and
that $F_1'$ be a geometrically trivial covering of~$F$.
\end{theorem}

The
% original p. 256
necessity of the condition being trivial, it remains to prove the
sufficiency. One easily reduces to the case where $S$ is
noetherian. Consider the Stein factorization $X\to T\to S$ of~$f$,
where $T$ is the spectrum of the Algebra $f_*(\cal{O}_X)$ over~$S$.
Since the fibers of~$X$ over~$S$ are geometrically connected, and $f$
is surjective, the morphism $T\to S$ is finite surjective and radicial,
hence \eqref{IX.4.10} every $T'$ étale over~$T$ comes by inverse image
from an $S'$ étale over~$S$. This reduces us to
proving~\Ref{IX.6.5}
with $S$ replaced in it by~$T$, \ie to the case where one assumes
$f_*(\cal{O}_X)=\cal{O}_S$. Consider then the Stein factorization
$X'\to S'\to S$ of the proper morphism $h\colon X'\to S$, where $S'$
is the spectrum of the Algebra $h_*(\cal{O}_{X'})$. The morphisms
$X'\to X$ and $X'\to S'$ define a canonical morphism
$$
X'\to X\times_S S',
$$
and our assertion is contained in the following one:

\begin{corollary}
\label{IX.6.6}
Let $f\colon X\to S$ be a proper morphism of locally noetherian
preschemes, such that $f_*(\cal{O}_X)=\cal{O}_S$, and let $X'$ be a
prescheme proper over~$X$. Consider the \emph{Stein} factorization
$X'\to S'\to S$ for $X'\to S$, and the canonical morphism $X'\to
X\times_S S'$. Let $s$ be a point of~$S$, $s'$ a point of~$S'$
over~$s$, corresponding to a connected component $F_1'$ of the fiber
$X'_s$ of~$X'$ at~$s$. In order that the morphism $X'\to X\times_S S'$
be an isomorphism over an open neighborhood $U'$ of~$s'$ étale
over~$S$, it is necessary and sufficient that $X'$ be étale over~$X$ at
the points of~$F_1'$, and that $F_1'$ be a geometrically trivial
covering of the fiber $F=X_s$.
\end{corollary}

The necessity again being trivial, it remains to prove the
sufficiency. The conclusion also means that a) the morphism deduced
from~$X'\to X\times_S S'$ by the change of base
$\Spec(\widehat{\cal{O}}_{s'})\to S'$ is an isomorphism, b) $S'$ is
étale over~$S$ at~$s'$, \ie $\widehat{\cal{O}}_{s'}$ is étale
over~$\widehat{\cal{O}}_s$. In this form, one sees that the conclusion
is invariant under the change of base $\Spec(\widehat{\cal{O}}_s)\to
S$. The hypotheses being likewise stable under this change of base, one
may therefore assume that $S$ is the spectrum of a complete noetherian
local ring. One may moreover obviously assume $X'$ connected, which
implies here that $S'=\Spec(\cal{O}_{s'})$, $F'=F_1'$.

Since
% original p. 257
the set of points of~$X'$ where $X'$ is étale over~$X$ is open and
contains the fiber $X'_{s'}=F'$, $X'$ being proper over~$S$, it
follows that $X'$ is étale over~$X$. Since it induces on~$F=X_s$ an
étale covering isomorphic to some $F\otimes_{\kres(s)}L$, where $L$ is
étale over~$\kres(s)$, it follows from~\Ref{IX.1.10} that it is
isomorphic to a covering of the form $X\times_S T$, with $T$ étale
over~$S$. (N.B. here again, it suffices to use that the functor
of~\Ref{IX.1.10} is fully faithful, which follows from the fact that a
formal isomorphism of coherent sheaves on~$X$ comes from an isomorphism
of these sheaves). Hence, if $T$ is defined by the algebra~$B$ finite
over~$A$, $X'$ is identified with the spectrum of the Algebra
$\cal{O}_X\otimes_A B$ over~$X$, from which it follows at once, since
$f_*(\cal{O}_X)=\cal{O}_S$, that $h_*(\cal{O}_{X'})$ is defined
by~$B$, hence the canonical homomorphism $X'\to X\times_S S'$ is none
other than the isomorphism under consideration $X'\isomto X\times_S T$.
This completes the proof.

\begin{corollary}
\label{IX.6.7}
Under the conditions of~\Ref{IX.6.5}, in order that there exist a
prescheme $S'$ étale over~$S$ and an $X$-isomorphism $X'\isomto
X\times_S S'$, it is necessary and sufficient that $X'$ be étale
over~$X$ and that for every $s\in S$, $X'_{s'}$ be a geometrically
trivial covering of~$X_s$.
\end{corollary}

Indeed, if this is so, $X'$ is the union of open sets $X_i'$ which are
isomorphic to inverse images of $S_i'$ étale over~$S$. One then easily
sees that these $S_i'$ glue into an $S'$ étale over~$S$, and that one
obtains an isomorphism $X'\isomto X\times_S S'$. One can say for
example that the $X_i'$ are endowed with descent data relative
to~$X\to S$, which necessarily glue into a descent datum on the whole
of~$X'$ relative to~$X\to S$. And since the latter is effective on
the~$X_i'$, it follows easily (thanks to a sorites forgotten in
\No\Ref{IX.4}) that it is effective. One can also state \Ref{IX.6.7}
in the following form:

\begin{corollary}
\label{IX.6.8}
Let $f\colon X\to S$ be a proper surjective morphism of finite
presentation, with geometrically connected fibers. Then $f$ is an
effective descent morphism for the fibered category of preschemes
finite étale over others. The functor $S'\mto X\times_S S'$ induces an
equivalence of the category of preschemes étale and finite over~$S$
with the category of preschemes
% original p. 258
étale and finite over~$X$ which induce on each fiber $X_s$ a
geometrically trivial covering.
\end{corollary}

\begin{remark}
\label{IX.6.9}
Let $f\colon X\to S$ be a proper and surjective morphism, with $Y$
locally noetherian; then $f$ factors into a morphism $X\to S'$
satisfying the hypothesis of~\Ref{IX.6.8}, and a finite surjective
morphism $S'\to S$ to which \Ref{IX.4.7} applies, hence into a product
of two morphisms which are \emph{universal effective descent
morphisms} for the fibered category of preschemes étale and finite
over others. One can conclude from this that $f$ itself is a universal
effective descent morphism for the fibered category under
consideration. One thus recovers \Ref{IX.4.12} by a different method.
\end{remark}

\begin{remark}
\label{IX.6.10}
The conclusion of~\Ref{IX.6.7} does not remain valid if one replaces
the hypothesis that $f$ is proper by: $X$ is of finite type over~$S$
and admits a section over~$S$ (so that $f$ is universally submersive
and a descent morphism for the fibered category of preschemes étale
over others), even when $S$ is the spectrum of a discrete valuation
ring and when $X'$ is an étale covering of~$X$. To see this, one
starts from a $Z$ proper over~$S$, whose general fiber
is a nonsingular rational curve, and whose special fiber $Z_0$
consists of two concurrent lines. For example, if $t$ is a uniformizer
of the valuation ring~$A$, one takes the closed subscheme $Z$
of~$\PP_A^2$ defined by the homogeneous equation $x^2+y^2+tz^2=0$
(homogeneous coordinates $x$, $y$, $z$). One takes for $X$ the
complement of the singular point $a$ of~$Z_0$ in the union $Z\cup
\PP_k^2$. The fibers of~$X$ are $\PP_k^1$ and $\PP_k^2 - a$, hence are
geometrically simply connected (\ie every étale covering of such a
fiber is geometrically trivial). However
% original p. 259
one easily constructs, proceeding as in \No \Ref{IX.4}, étale
coverings of~$X$ which do not come from étale coverings of~$S$, by
gluing trivial coverings of~$Z-a$ and of~$\PP_k^2-a$. It is possible
on the other hand that the conclusion of~\Ref{IX.6.7} persists if one
replaces in it the hypothesis of properness by that $X$ be
\emph{universally open} of finite presentation
over~$S$\kern1pt\footnote{This is now proved, $g$ being only
universally open and surjective; \cf SGA~4 XV~1.15.}. This is true at
least if one assumes that the fibers of~$X$ over~$S$ are geometrically
irreducible, and not only geometrically connected. Let us only point
out that in this question one can reduce to the case where $S$ is the
spectrum of a complete discrete valuation ring with algebraically
closed residue field.
\end{remark}

The interpretation of~\Ref{IX.6.7} in terms of the fundamental group is
the following:

\begin{corollary}
\label{IX.6.11}
Let $f\colon X\to S$ be a proper surjective morphism of finite
presentation, with geometrically connected fibers. Suppose $X$, hence
$S$, connected. Let $a$ be a geometric point of~$X$, $b$ its image
in~$S$, and for every $s\in S$, let us choose an algebraic closure
$\overline{\kres(s)}$ of $\kres(s)$, a geometric point $a_s$ of~$S_s$
with values in this extension, and a class of paths from~$a_s$ to~$a$,
whence a homomorphism $\pi_1(\overline X_s,a_s)\to \pi_1(X,a)$, where
$\overline X_s=X_s\otimes_{\kres(s)}\overline{\kres(s)}$. Then the
homomorphism $\pi_1(X,a)\to \pi_1(S,b)$ is surjective, and its kernel
is the closed invariant subgroup of~$\pi_1(X,a)$ generated by the
images of the $\pi_1(\overline X_s,a_s)$.
\end{corollary}

\begin{remark}
\label{IX.6.12}
Under the conditions of~\Ref{IX.6.7}, assuming $S$ noetherian, one
easily sees that the set of points $s\in S$ such that $S'_{s'}$ is
geometrically trivial over~$X_s$ is a constructible set; if $X'$ is
proper over~$X$, it is even open, as one sees from~\Ref{IX.6.6}. This
therefore allows one, if $S$ is a Jacobson prescheme (for example of
finite type over a field), or
% original p. 260
$X'$ is proper over~$X$, to restrict oneself, in order to verify the
conditions of~\Ref{IX.6.7}, to the points $s$ of~$S$ which are
\emph{closed}. Likewise, in~\Ref{IX.6.11} it then suffices to take the
$\pi_1(\overline X_s,a_s)$ for the points of~$S$ which are closed.
\end{remark}

\begin{thebibliography}{D}{IX.7}

\bibitem[D]{IX.D} J\ptbl \textsc{Giraud}, M\'ethode de la
descente,
M\'emoire \no 2
de la Soc.\ Math.\ Fr.,~1964.

\bibitem[1]{IX.1} A\ptbl \textsc{Grothendieck}, G\'eom\'etrie Alg\'ebrique et
G\'eom\'etrie Formelle, S\'eminaire Bourbaki t\ptbl 11, 1959,
\No 182.
\end{thebibliography}
